\documentclass[11pt]{article}

% --- Geometry & typography -------------------------------------------------
\usepackage[a4paper,margin=1in]{geometry}
\usepackage[T1]{fontenc}
\usepackage[utf8]{inputenc}
\usepackage{lmodern}
\usepackage{microtype}
\usepackage{textcomp}

% --- Math & tables ---------------------------------------------------------
\usepackage{amsmath,amssymb}
\usepackage{booktabs}
\usepackage{array}
\usepackage{multirow}
\usepackage{tabularx}
\usepackage{siunitx}
\sisetup{detect-all,table-number-alignment=center}

% --- Graphics & cross-refs -------------------------------------------------
\usepackage{graphicx}
\graphicspath{{../../results/figures/}}
\usepackage[font=small,labelfont=bf,labelsep=period]{caption}
\usepackage[dvipsnames]{xcolor}
\definecolor{linkblue}{RGB}{30,80,150}
\usepackage[colorlinks=true,linkcolor=linkblue,citecolor=linkblue,urlcolor=linkblue]{hyperref}
\usepackage{enumitem}
\setlist[itemize]{leftmargin=1.2em,itemsep=0.2em,topsep=0.3em}

% --- Headers ---------------------------------------------------------------
\usepackage{authblk}
\renewcommand\Authsep{\quad}
\renewcommand\Authands{\quad}

% --- Table helpers ---------------------------------------------------------
\newcommand{\dmse}{$\Delta$ input MSE}
\newcommand{\dauc}{$\Delta$ AUC}

\title{\textbf{Mechanistic Decomposition of BrainOmni's Tokenizer Source Basis:\\
Causal Necessity, Diagnostic Sufficiency, and Modality-Dependent Asymmetries}}
\author{Anonymous}
\date{Draft, May~2026}

\begin{document}
\maketitle

\begin{abstract}
\noindent
BrainOmni is a brain foundation model whose front end --- the BrainTokenizer ---
quantizes raw EEG/MEG signals through a 16-source residual vector-quantization (RVQ)
bottleneck. We ask what each of the 16 \emph{source latents} encodes and whether the
directions inside them carry interpretable, causally important, and diagnostic-relevant
structure. We train sparse autoencoders (SAEs) on the tokenizer's pre-quantizer
activations and identify, for each source, a single \emph{source-prototype direction}
in the standardized activation space.  Across two datasets (PD31 EEG, 10~subjects,
188~segments; MEG-MMI, 6~subjects, 159~segments) and four SAE seeds we show that
(i)~the source-prototype direction is seed-invariant (pairwise cosine
$\geq 0.998$ in PD31, $\geq 0.997$ in MEG);
(ii)~zero-ing the prototype projection from source-conditioned rows damages tokenizer
reconstruction by 0.012--0.020 in input MSE while removing a random unit direction
does not;
(iii)~on PD31, a linear classifier built from the \emph{single} prototype 1-d
projection matches the full 256-d activation classifier (stratified-5-fold
AUC~$0.760$ vs $0.748$), and on MEG it captures 87.7\,\% of the
leave-one-subject-out baseline accuracy;
(iv)~cross-subject patching of the prototype-orthogonal residual transports
$\sim 21$--$8\times$ more donor-specific structure than patching the projection,
consistently across both modalities;
(v)~on MEG, sources~2 and~9 dissociate cleanly --- source~9 is more
reconstruction-causal but transports zero donor structure and barely shifts the
diagnostic classifier, while source~2 carries label-aligned, subject-specific
transferable structure; and
(vi)~the source-3 PD31 residual decomposes into $\sim 44\,\%$ label-aligned and
$\sim 87\,\%$ label-orthogonal donor delta.
The BrainTokenizer's 16-source basis is therefore decomposing brain signals into
qualitatively distinct latents --- not a single 1-d diagnostic axis --- and the
prototype direction is the strongest single-axis interpretability handle in the model.
\end{abstract}

% =========================================================================
\section{Introduction}\label{sec:intro}
% =========================================================================

Brain foundation models are starting to scale to multiple modalities (EEG, MEG,
BIDS-formatted data) and downstream tasks (sleep staging, depression detection,
Parkinson's classification). Following the success of self-supervised
pre-training on raw {EEG} \cite{bendr2021,banville2021} and dedicated foundation
models for non-invasive \cite{labram2024,biot2023} and invasive
\cite{brant2023,brant2_2024} brain signals, as well as for {fMRI}
\cite{brainlm2024} and audio-aligned {MEG} \cite{defossez2023brain},
BrainOmni~\cite{brainomni2025} is one such model unifying {EEG} and {MEG}; its
front end is the \textbf{BrainTokenizer}, a residual vector-quantizer
\cite{vqvae2017,vqvae2_2019,soundstream2021,encodec2022} that compresses
arbitrary brain recordings into a discrete code sequence by routing each input
row through 16 learned \emph{source latents} with shared codebooks. While
downstream BrainOmni reports impressive transfer numbers, the question of
\emph{what each source latent encodes} --- and whether they decompose brain
signals into interpretable structure or simply hash the input --- has not been
investigated.

This paper makes that question concrete. We treat each tokenizer source as a
candidate ``interpretable latent'' and ask, for the two most-causal sources on
each of two datasets, whether there is a low-dimensional direction inside the
source's activation distribution that:
\begin{enumerate}[leftmargin=1.4em,itemsep=0.2em,topsep=0.3em]
\item \textbf{Replicates across SAE seeds.} Different sparse-autoencoder fits
should converge on the same direction.
\item \textbf{Is causally necessary} for tokenizer reconstruction. Zero-ing it
should damage reconstruction MSE substantially more than zero-ing a random
direction.
\item \textbf{Is diagnostic-aligned.} The 1-d projection should carry information
that a downstream linear classifier uses to discriminate diagnostic labels
(HC vs PD; HV vs MDD).
\item \textbf{Survives cross-subject transport.} Patching it from a donor segment
to a recipient segment should move the recipient's reconstruction toward the
donor's pattern.
\item \textbf{Compares favorably to standard linear baselines} (PCA, FastICA on
the same activations).
\end{enumerate}

We find a single direction satisfying (1)--(2) on both datasets --- the
\emph{source-prototype direction}, computed from a trained SAE's encoding of
source-conditioned rows. We then probe the geometry of this direction along the
four remaining axes. Our headline findings:
\begin{itemize}
\item \textbf{The source-prototype direction is seed-invariant} to
$<\!0.003$ cosine units across 4 SAE seeds in both PD31 EEG and MEG-MMI.
\item \textbf{The 1-d prototype projection is sufficient for diagnostic
classification on PD31:} stratified-5-fold AUC~$=0.760$ from the
\emph{projection only} matches the full 256-d activation baseline ($0.748$).
Removing the prototype direction barely changes the classifier
($\Delta\text{AUC}=-0.007$); removing a random unit direction is
bit-identical to baseline.
\item \textbf{Cross-subject residual patching transports $8\times$ more donor
structure than projection patching} on PD31, and $21\times$ more on MEG.
Donor-specific source variation lives mainly in the \emph{prototype-orthogonal}
residual, while the prototype direction itself is approximately shared across
subjects.
\item \textbf{PCA and ICA invert their roles between EEG and MEG.} On PD31
source~3, FastICA on the standardized activation rows finds a direction with
cosine $0.71$ to the SAE prototype and damages reconstruction \emph{more} than
the SAE prototype itself; PCA's most-aligned component damages negligibly.
On MEG source~2 the picture inverts: PCA finds a direction with cosine $0.72$
that damages reconstruction $2\times$ the SAE prototype, while ICA's best
component sits at cosine $0.19$ (2-subject) or $0.75$ (6-subject expanded).
\item \textbf{Cross-source dissociation on MEG.} Source~9 is more
reconstruction-causal than source~2 ($\Delta$MSE $0.023$ vs $0.020$) but
transports \emph{zero} donor-specific structure under patching (residual
cosine $-0.004$) and barely shifts the diagnostic classifier
($\Delta\text{AUC}=-0.013$ vs $-0.047$ for source~2). The two MEG sources
factorize qualitatively different regularities.
\end{itemize}

We end with a refined account of the residual structure on PD31: decomposing the
prototype-orthogonal residual into label-aligned and label-orthogonal components,
$\sim 44\,\%$ of the residual donor delta lives in the 1-d HC$\leftrightarrow$PD
axis but $\sim 87\,\%$ lives in the orthogonal complement (mostly
subject-specific noise). The ``label-structured residual'' claim from the
project's earlier reports holds, but the bulk of the residual donor transport is
label-\emph{orthogonal}.

Together these results give a six-property fingerprint of the BrainTokenizer's
source-prototype directions: seed-stable, reconstruction-necessary,
diagnostic-sufficient (on PD31), cross-subject portable through the residual,
modality-dependent in their relation to PCA/ICA baselines, and qualitatively
distinct across sources within MEG.

% =========================================================================
\section{Background}\label{sec:background}
% =========================================================================

\subsection{BrainOmni and the BrainTokenizer}

BrainOmni is a brain foundation model with two stages: a frozen
\textbf{BrainTokenizer} (5.05~M parameters) that maps any continuous EEG/MEG
recording to a discrete token sequence, and a \textbf{BrainOmni encoder}
(13.5~M parameters in \texttt{tiny}, larger in \texttt{base}) that processes
those tokens for downstream tasks. The tokenizer is a residual vector-quantizer
with 16 \emph{source latents} and 4 quantization levels per source (codebook
size 512 each). Its forward pass ingests $(x, \mathrm{pos},
\mathrm{sensor\_type})$ triples --- a multichannel signal $x$, sensor positions
$\mathrm{pos}$, and sensor types $\mathrm{sensor\_type}$ --- and produces, per
spatiotemporal patch, a $16\!\times\!4\!\times\!256$-dimensional residual
representation that is then decoded back to the input space.

The 16 source latents are not specialized to a single sensor or frequency band
by construction: each source can attend to any subset of the input channels
through shared sensor embeddings. Whether the model uses each source for a
distinct purpose, or hashes inputs uniformly across sources, is the empirical
question we address.

\subsection{Datasets}

We use two preprocessed OpenNeuro subsets:
\begin{itemize}
\item \textbf{PD31 EEG} (OpenNeuro \texttt{ds002778-1.0.5})
\cite{ds002778}: 10~subjects (5 healthy controls / 5 Parkinson's),
32-channel EEG resting-state, 188 30-second segments after
preprocessing. The ``expanded10'' subset; the original ``4-subject''
subset (76~segments) is also referenced where multi-seed analyses
preceded the expansion.
\item \textbf{MEG-MMI} (OpenNeuro \texttt{ds003568-1.0.4})
\cite{ds003568}: 6~subjects (3~healthy volunteers / 3~major
depressive disorder), 269-channel MEG monetary-gambling
mood-induction task in the spirit of the MID paradigm
\cite{knutson2001mid}, 159 30-second segments after preprocessing.
The ``expanded6'' subset, expanded from an original 2-subject set
(20~segments) used for early multi-seed work.
\end{itemize}
Both use the upstream BrainOmni preprocessing pipeline (band-pass + notch
filter, resample to 256~Hz, 30-second windowing). All MEG analyses use a
10-second clip per segment for source-prototype interventions to keep
per-segment compute tractable; uncropped (full 30-second) analyses are
reported separately.

\subsection{Sparse autoencoders on tokenizer activations}

Sparse autoencoders are an interpretability tool with a long lineage in
neuroscience-inspired sparse coding \cite{olshausen1996} and have, in
recent work, been used to decompose language-model activations into
interpretable directions
\cite{bricken2023monosemanticity,cunningham2024sae,templeton2024scaling,gao2024scaling,rajamanoharan2024gated,lieberum2024gemma}.
For each dataset we exported a tokenizer activation bundle
(rows $\times$ 256-d) at the \texttt{encoder\_pre\_quantizer} hook (the
standardized activation entering the residual quantizer) and trained a
top-$k$ SAE \cite{gao2024scaling} with hidden dimension 1024, $k = 64$,
and 10 epochs. Quality metrics: PD31 expanded10 SAE val MSE $0.0294$ /
explained variance $0.971$; MEG-MMI expanded6 SAE val MSE $0.0236$ /
explained variance $0.976$. Across 4 SAE seeds, the same architecture
and dataset reproduce within 5\,\% on these metrics.

% =========================================================================
\section{The source-prototype direction}\label{sec:prototype}
% =========================================================================

For each source latent $s$ we want a single direction $p_s$ in the SAE's 256-d
standardized activation space that summarizes ``what source $s$ is doing on
average''. We compute it as follows. Let $R_s$ be the matrix of standardized
SAE input rows from source $s$ and $E(\cdot)$ the SAE encoder. Then
\begin{align*}
\mathrm{feature\_mean}_s &= \mathrm{mean\ over\ rows\ of\ } E(R_s) \quad (\text{1024-d hidden})\\
\mathrm{prototype}_s     &= \mathrm{feature\_mean}_s \cdot W_{\mathrm{dec}}^\top \quad (\text{256-d})\\
\mathrm{prototype}_s     &\leftarrow \mathrm{prototype}_s / \|\mathrm{prototype}_s\|
\end{align*}
where $W_{\mathrm{dec}}$ is the SAE decoder's row-normalized weight matrix.
This is the \emph{mean of decoded source-$s$ features}, projected back to
activation space and unit-normalized. The construction treats the
prototype as a single direction in the standardized activation space,
in the spirit of the linear-representation hypothesis
\cite{park2024linear} and prior dictionary-learning approaches to feature
discovery \cite{olshausen1996,bell1995infomax}.

Two practical notes. First, the SAE is trained jointly across sources and
quantizers --- it does not see source identity at training time. The prototype
direction is therefore a \emph{learned summary} of the source's typical
activation pattern, not a hand-coded mask. Second, the prototype is computed
from the SAE's decoder rather than its encoder so that interventions on
activation space correspond to a clean linear direction; encoder-side
directions would entangle the top-$k$ sparsity decision \cite{gao2024scaling}.

\subsection{Seed stability of the prototype direction}

We trained 4 SAEs at seeds 0--3 on each dataset with otherwise identical config,
then computed pairwise cosine similarities between the corresponding
source-prototype directions. Figure~\ref{fig:seed-stability} shows the result:
pairwise cosine $\geq 0.998$ in PD31 source~3 and $\geq 0.997$ in MEG source~2
across all 6 same-dataset seed pairs. The prototype direction is essentially
seed-invariant. By comparison, individual SAE \emph{features} in the same
models have mean active-mutual cosine $\approx 0.5$ across seeds --- feature
identity is not stable, but the source-level prototype is.

This stability also persists to the expanded MEG SAE: the 4-seed expanded6
sweep gives source-2 \dmse{} $= 0.01988 \pm 0.00054$ (std 2.7\,\% of mean).
Source~9 is stable at \dmse{} $= 0.0248 \pm 0.0013$ across seeds 0--2, with
seed~3 a clean outlier at $0.00078$ --- same prototype magnitude (mean abs
projection $11.02$ vs $11.21$--$11.32$) but oriented away from the
input-reconstruction error axis. This is the only seed-instability of the
entire sweep and is worth flagging in any paper extension.

\begin{figure}[t]
\centering
\includegraphics[width=0.92\linewidth]{fig_seed_stability.pdf}
\caption{\textbf{Pairwise cosine of the source-prototype direction across 4 SAE
seeds.} Off-diagonal cosine is $\geq 0.998$ on PD31 source~3 (left) and
$\geq 0.997$ on MEG source~2 (right). The source-level prototype is
essentially seed-invariant, in contrast with the individual SAE features it
is built from (mean active-mutual cosine $\approx 0.5$).}
\label{fig:seed-stability}
\end{figure}

% =========================================================================
\section{Reconstruction necessity}\label{sec:necessity}
% =========================================================================

For each segment with source-$s$ rows in the standardized activation space, we
replace each row $r$ with $r - (r \cdot p_s)\,p_s$ and re-run the quantizer +
decoder forward, in the activation-patching style introduced by
\cite{vig2020causal,rome2022,wang2023ioi,conmy2023acdc} and unified
under causal abstraction \cite{geiger2024abstraction}. We measure
$\Delta\mathrm{input\,MSE} = \mathrm{MSE}(\hat x_{\mathrm{intervened}}, x) -
\mathrm{MSE}(\hat x_{\mathrm{baseline}}, x)$.
Averaged across 76 segments (PD31 4-subject), 188 segments (PD31 expanded10),
and 40 balanced segments (MEG-MMI expanded6):

\begin{table}[h]
\centering
\small
\begin{tabular}{lccrl}
\toprule
\textbf{dataset} & \textbf{source} & \textbf{seeds} & \textbf{\dmse{}} & \textbf{random-dir.\ \dmse{}} \\
\midrule
PD31 (4 subj, 76 seg)     & src 3 & 4 & $0.01190 \pm 0.00015$ & $0.000$ (no effect) \\
PD31 (10 subj, 40 bal.)   & src 3 & 1 & $0.01164$             & --- \\
MEG (2 subj, 20 seg)      & src 2 & 4 & $0.01174 \pm 0.00062$ & $0.000$ \\
MEG (6 subj, 40 bal.)     & src 2 & 4 & $0.01988 \pm 0.00054$ & $0.000$ \\
MEG (6 subj, 40 bal.)     & src 9 & 3+1 & $0.0248 \pm 0.0013$ (seed 3: $0.00078$) & $0.000$ \\
\bottomrule
\end{tabular}
\caption{Mean $\Delta$ input MSE when the source-prototype projection is
zero-ed from source-conditioned rows. Random-direction removal is a no-op.
\label{tab:reconstruction}}
\end{table}

Three observations. \emph{(i)}~Prototype-removal damage is much larger than
random-direction removal --- confirming the prototype direction is not just any
direction in 256-d. \emph{(ii)}~The damage is $\sim$constant across SAE seeds
and across the 4-to-10-subject PD31 expansion ($0.01190 \to 0.01164$).
\emph{(iii)}~On MEG the damage \emph{increases} with the 2 $\to$ 6 subject
expansion ($0.01174 \to 0.01988$, a $1.7\times$ increase) --- the original
2-subject result \emph{under}-estimated the prototype's causal weight, not over.

\begin{figure}[h]
\centering
\includegraphics[width=0.78\linewidth]{fig_prototype_damage.pdf}
\caption{\textbf{Source-prototype removal damage is seed-stable across 4
seeds.} Mean $\Delta$ input MSE on each SAE seed for PD31 source~3 (4-subject
76-segment set, blue) and MEG source~2 (2-subject 20-segment set, orange).
Three of four seeds match within 0.0005 in PD31 and within 0.0014 in MEG;
seed~3 shows a single mild MEG outlier ($0.0107$).}
\label{fig:prototype-damage}
\end{figure}

% =========================================================================
\section{The 1-d prototype projection is diagnostically sufficient on PD31}\label{sec:probe}
% =========================================================================

The reconstruction-necessity result tells us the prototype direction matters
for the tokenizer's input fidelity. We now ask whether it carries diagnostic
information --- whether a linear classifier can use it to discriminate HC vs
PD on PD31 and HV vs MDD on MEG-MMI.

\subsection{Probe design}

For each segment we compute a feature vector by averaging the standardized
\texttt{encoder\_pre\_quantizer} rows across all 16 sources (256-d feature).
We use linear classifier probes \cite{alain2017probes} as the
diagnostic readout, with the standard probing-classifier caveats around
selectivity and confound \cite{belinkov2022probing,hewitt2019probes}.
We optionally apply one of three interventions to the source-$s$ rows
before pooling:
\begin{itemize}
\item \textbf{Baseline}: no intervention.
\item \textbf{Prototype removed}: replace each row with $r - (r \cdot p_s)p_s$.
\item \textbf{Prototype only}: replace each row with $(r \cdot p_s)\,p_s$
(1-d projection, rank-1).
\item \textbf{Random direction removed}: same as ``prototype removed'' but
along a fixed random unit direction.
\end{itemize}
We train an L2 logistic regression sweeping
$C \in \{10^{-3}, 10^{-2}, 0.1, 1, 10\}$ under two CVs:
leave-one-subject-out (LOSO; cross-subject generalization) and stratified
5-fold (within-subject leakage allowed; upper bound on label structure in the
activations).

\subsection{PD31 expanded10 results}

\begin{table}[h]
\centering
\small
\begin{tabular}{lcccc}
\toprule
\textbf{setting} & \textbf{strat.\ 5-fold acc.} & \textbf{strat.\ 5-fold AUC} & \textbf{LOSO acc.} & \textbf{LOSO AUC} \\
\midrule
baseline                 & $0.681$ & $0.748$           & $0.277$ & $0.177$ \\
prototype removed        & $0.676$ ($-0.005$) & $0.741$ ($-0.007$) & $0.266$ & $0.158$ \\
\textbf{prototype only}  & $\mathbf{0.686}$ ($+0.005$) & $\mathbf{0.760}$ ($+0.012$) & $0.266$ & $0.154$ \\
random dir.\ removed     & $0.681$ ($+0.000$) & $0.749$ ($+0.001$) & $0.277$ & $0.177$ \\
\bottomrule
\end{tabular}
\caption{HC vs PD linear-probe metrics on PD31 expanded10. The 1-d prototype
projection (rank-1) matches the full 256-d activation baseline.
\label{tab:pd31-probe}}
\end{table}

(See Figure~\ref{fig:downstream-probe}, left panel.) Three clean conclusions:
\begin{enumerate}[leftmargin=1.4em,itemsep=0.2em,topsep=0.3em]
\item \textbf{The prototype direction is sufficient.} Replacing the full
256-d activation with its 1-d projection onto the prototype direction
(zero-ing the orthogonal complement) keeps classifier performance ---
$0.686$ / $0.760$ vs the full-baseline $0.681$ / $0.748$. The 1-d direction
encodes essentially all of the diagnostic signal that the linear classifier
uses.
\item \textbf{The prototype direction is not necessary.} Removing the
prototype direction degrades the classifier by less than 1\,\%. The
diagnostic signal is \emph{redundantly} duplicated between the prototype
direction and the prototype-orthogonal residual.
\item \textbf{Removing a random direction is a no-op.} $\Delta$Accuracy
$=0.000$, $\Delta$AUC $=+0.001$ --- confirming the prototype-removed result
is not a generic ``remove any 1-d direction'' effect.
\end{enumerate}
LOSO collapses to below-chance for all settings because with 10 subjects the
per-subject activation bias overwhelms the label signal under a strict
cross-subject split (each fold has only one held-out subject contributing
9\,\% of the segments).

\subsection{4-seed robustness on the original 4-subject set}

Same probe applied to each of the 4 PD31 SAE seeds trained on the original
4-subject 76-segment set (2 HC + 2 PD subjects). Stratified 5-fold mean
$\pm$ std across the 4 seeds: baseline $0.895 \pm 0.000$ acc / $0.931 \pm 0.000$
AUC; prototype removed $0.891 \pm 0.007$ / $0.942 \pm 0.001$; prototype only
$0.852 \pm 0.007$ / $0.926 \pm 0.001$; random dir.\ removed $0.895 \pm 0.000$ /
$0.935 \pm 0.000$. All four seeds agree to within $0.001$ AUC; the 1-d
sufficiency claim is seed-robust at smaller sample size as well. (The higher
absolute scores reflect 2-subject confounds, not better generalization.)

\subsection{MEG expanded6 results}

LOSO becomes meaningful at MEG ($N\!=\!6$ subjects, $\sim\!26$ segments per
subject):

\begin{table}[h]
\centering
\small
\begin{tabular}{lcccc}
\toprule
\textbf{setting} & \textbf{LOSO acc.} & \textbf{LOSO AUC} & \textbf{\dauc} & \\
\midrule
baseline                       & $0.730$ & $0.715$ & ---       \\
prototype removed              & $0.698$ & $0.710$ & $-0.005$  \\
\textbf{prototype only (src 2)} & $0.654$ & $0.668$ & $\mathbf{-0.047}$ \\
\textbf{prototype only (src 9)} & $0.723$ & $0.701$ & $\mathbf{-0.013}$ \\
random dir.\ removed           & $0.730$ & $0.715$ & $0.000$   \\
\bottomrule
\end{tabular}
\caption{HV vs MDD linear-probe metrics on MEG expanded6 under LOSO.
\label{tab:meg-probe}}
\end{table}

(See Figure~\ref{fig:downstream-probe}, right panel.) On MEG the 1-d source-2
projection captures 87.7\,\% of baseline accuracy and 93\,\% of baseline AUC
under genuine cross-subject generalization. Source~9, by contrast, captures
only 99\,\% of baseline accuracy and 98\,\% of baseline AUC --- i.e.\ it is
weakly diagnostic. This source dissociation is taken up in
Section~\ref{sec:dissociation}.

\begin{figure}[t]
\centering
\includegraphics[width=0.95\linewidth]{fig_downstream_probe.pdf}
\caption{\textbf{Downstream classifier under prototype-direction
interventions, both modalities.} Stratified-5-fold accuracy (left) and AUC
(right) for the two diagnostic-label tasks. The PD31 (orange) prototype-only
column matches the baseline column, demonstrating 1-d sufficiency. The MEG
(blue) classifier is dominated by within-subject structure and is largely
robust to all four interventions at this CV; the genuine cross-subject
collapse appears under LOSO (Table~\ref{tab:meg-probe}). Dashed line at
0.5~marks chance.}
\label{fig:downstream-probe}
\end{figure}

% =========================================================================
\section{Cross-subject patching: residual versus projection}\label{sec:patching}
% =========================================================================

If the prototype direction is \emph{causally necessary} and
\emph{diagnostic-sufficient}, what does the \emph{residual} (orthogonal
complement) carry? We test cross-subject portability: pick 10 cross-subject
donor/recipient segment pairs, replace the recipient's source-$s$ rows with
the donor's only along one of three components --- prototype projection,
prototype-orthogonal residual, or full row --- and measure the cosine
alignment between the patched delta and the donor delta in residual space.

\begin{table}[h]
\centering
\small
\begin{tabular}{llrr}
\toprule
\textbf{dataset} & \textbf{component} & \textbf{donor-align.\ cos.} & \textbf{mean \dmse{}} \\
\midrule
PD31 expanded10        & residual    & $\mathbf{0.391}$ & $+0.0273$ \\
PD31 expanded10        & projection  & $0.048$          & $-0.0032$ \\
MEG expanded6 src 2    & residual    & $\mathbf{0.168}$ & $+0.033$  \\
MEG expanded6 src 2    & projection  & $0.008$          & $-0.001$  \\
MEG expanded6 src 9    & residual    & $\mathbf{-0.004}$ & $+0.0205$ \\
MEG expanded6 src 9    & projection  & $-0.008$          & $+0.0003$ \\
\bottomrule
\end{tabular}
\caption{Cross-subject patching: donor-direction alignment cosine and mean
$\Delta$ input MSE for projection-only vs residual-only patching, three
sources across two datasets (10 cross-subject pairs each).
\label{tab:patching}}
\end{table}

Three results:
\begin{enumerate}[leftmargin=1.4em,itemsep=0.2em,topsep=0.3em]
\item \textbf{The residual is the carrier of donor-specific structure for
src-3 EEG and src-2 MEG.} Residual donor alignment is $8\times$ larger than
projection donor alignment on PD31 and $21\times$ larger on MEG src~2.
\item \textbf{The MEG src-9 residual carries zero transferable donor
structure.} Both residual and projection patching cosines are
indistinguishable from zero. Source~9 looks like a shared population
baseline rather than a subject-specific residual carrier.
\item \textbf{The honest 6-subject MEG residual donor-alignment is $0.168$,
not $0.99$.} The original 2-subject MEG residual cosine of $\sim 0.99$ was an
artifact of having only one HV/MDD pair worth of contrast; with 6 subjects,
cross-subject diversity is real.
\end{enumerate}

\begin{figure}[t]
\centering
\includegraphics[width=0.85\linewidth]{fig_patch_decomposition.pdf}
\caption{\textbf{Donor-specific source variation lives in the
prototype-orthogonal residual.} Mean cosine alignment between the patched
delta and the donor delta across 10 cross-subject pairs, projection-only vs
residual-only patching, three sources across two datasets. The MEG src-9
residual is a clean zero --- source 9 transports no donor-specific
structure, in striking contrast to the residuals of MEG src 2 and PD31
src 3.}
\label{fig:patch-decomposition}
\end{figure}

\subsection{Same- vs different-label residual transport on PD31}

To test whether residual transport carries \emph{label} information, we built
similarity-matched 10-pair sets --- 10 same-label cross-subject pairs
(HC$\leftrightarrow$HC, PD$\leftrightarrow$PD) and 10 different-label
cross-subject pairs (HC$\leftrightarrow$PD), with mean baseline residual cosine
matched to within $0.0025$. Under residual patching:

\begin{table}[h]
\centering
\small
\begin{tabular}{lcc}
\toprule
\textbf{pair set} & \textbf{seed-0 donor align.} & \textbf{seed-1 donor align.} \\
\midrule
same label      & $\mathbf{0.4032}$ & $\mathbf{0.4033}$ \\
different label & $0.3436$          & $0.3438$          \\
\bottomrule
\end{tabular}
\caption{Same- vs different-label residual transport on PD31 source~3. Two
SAE seeds; 10 similarity-matched pairs each.\label{tab:label-transport}}
\end{table}

The same-label transport advantage is real, identical across two SAE seeds,
and survives the similarity matching. The PD31 source-3 residual carries
cross-subject, label-aware structure: same-label pairs transport more than
different-label pairs of equal baseline similarity.

\begin{figure}[t]
\centering
\includegraphics[width=0.95\linewidth]{fig_label_transport.pdf}
\caption{\textbf{Same-label residuals transport more, at matched intervention
magnitude.} Left: donor-direction alignment of same-label vs different-label
residual transport on PD31 source~3, two SAE seeds, 10 similarity-matched
pairs each. Right: $\Delta$ input MSE in the same runs (same-label and
different-label interventions are within $0.001$ MSE of each other --- the
alignment difference is not driven by a magnitude mismatch).}
\label{fig:label-transport}
\end{figure}

\subsection{Decomposing the PD31 residual}

Where in the residual does the label information live? We decomposed the
source-3 prototype-orthogonal residual into: full residual, subject mean
(per-subject mean residual), subject-specific (segment $-$ subject mean),
label-aligned (1-d projection onto HC$\leftrightarrow$PD axis), and
label-orthogonal (residual $-$ label-aligned). For 10 random cross-label
pairs we measured the cosine alignment between the patched component delta
and the donor delta:

\begin{table}[h]
\centering
\small
\begin{tabular}{lc}
\toprule
\textbf{component} & \textbf{cosine to donor delta} \\
\midrule
full residual           & $1.000$ (by construction) \\
label-orthogonal        & $\mathbf{0.867}$ \\
subject-specific        & $0.776$ \\
subject mean            & $0.592$ \\
label-aligned (1-d)     & $\mathbf{0.437}$ \\
\bottomrule
\end{tabular}
\caption{Component-wise donor alignment of the PD31 source-3
prototype-orthogonal residual under cross-subject patching.
\label{tab:residual-decomp}}
\end{table}

So the source-3 residual donor delta decomposes as $\sim 44\,\%$ label-aligned
+ $\sim 87\,\%$ label-orthogonal (subject-specific noise dominates, but the
1-d HC$\leftrightarrow$PD axis carries a meaningful $44\,\%$ share). The two
components partially overlap because the residual basis is not orthogonal to
the label axis. HC vs PD label-mean cosine is $0.821$ --- the two label-mean
residuals are similar, so the label axis is a small but real correction on top
of a large shared mean.

% =========================================================================
\section{Cross-source dissociation on MEG: src 2 vs src 9}\label{sec:dissociation}
% =========================================================================

Combining Sections~\ref{sec:necessity}--\ref{sec:patching} yields a clean
dissociation between two MEG source latents. The original 2-subject MEG study
fixated on source~2 because it had the largest ablation footprint at small
sample size. With 6 subjects and uncropped 30-second segments, source ablation
reveals that \textbf{source~9 is the most-causal MEG source}
($\Delta$MSE $0.0503$ vs source~2's $0.0396$), and that the dominant
ablation-damage sensor cluster shifts from $\{244, 115, 251, 122, 245\}$
(small sample) to $\{185, 237, 244, 108, 156\}$ (uncropped 6-subject).

We then ran the full mechanistic pipeline on source~9 using the same
expanded6 SAE. The contrast with source~2:

\begin{table}[h]
\centering
\small
\begin{tabular}{lccc}
\toprule
\textbf{property} & \textbf{src 2} & \textbf{src 9} & \textbf{ratio} \\
\midrule
recon.-causal \dmse{}                     & $0.0197$ & $0.0232$ & $1.18\times$ \\
diagnostic LOSO \dauc{} (prototype-only)  & $-0.047$ & $-0.013$ & $0.28\times$ \\
residual donor-align.\ cos.\ (10 pairs)   & $+0.168$ & $-0.004$ & $\sim$zero \\
projection donor-align.\ cos.             & $+0.008$ & $-0.008$ & $\sim$zero \\
\bottomrule
\end{tabular}
\caption{Source~2 vs source~9 dissociation on MEG-MMI expanded6.
\label{tab:dissociation}}
\end{table}

So:
\begin{itemize}
\item \textbf{Source~2 is ``subject-specific deviation''}: less reconstruction
weight, more diagnostic weight, residual carries cross-subject transferable
structure.
\item \textbf{Source~9 is ``shared population baseline''}: more reconstruction
weight, less diagnostic weight, residual transports nothing.
\end{itemize}

The two sources factorize qualitatively different MEG regularities. The
BrainTokenizer's source basis is therefore not a single 1-d diagnostic axis
dressed up as 16 sources; it is a real factorization with at least two
distinct mechanistic roles. This is the strongest single argument in this
paper that the source latents are a meaningful mechanistic decomposition rather
than an arbitrary hash.

% =========================================================================
\section{PCA and FastICA baselines}\label{sec:baseline}
% =========================================================================

The SAE is one decomposition method among several. We compared its
source-prototype direction to two distribution-level decompositions:
{PCA}(64) and {FastICA}(64)
\cite{bell1995infomax,hyvarinen2000,makeig1996} on the same standardized
source-conditioned activation rows.
For each method we picked the component with the largest absolute cosine to the
SAE prototype direction, then ran the same \texttt{remove\_projection}
intervention along that single direction.

\subsection{PD31 source~3 (76 segments, seed 0 SAE)}

\begin{table}[h]
\centering
\small
\begin{tabular}{lcc}
\toprule
\textbf{method} & \textbf{cos.\ to SAE prototype} & \textbf{mean \dmse{}} \\
\midrule
SAE prototype                  & $1.000$  & $+0.0117$ \\
best PCA component (PC4)       & $-0.716$ & $-0.0013$ \\
best FastICA component (IC64)  & $-0.714$ & $\mathbf{+0.0172}$ \\
\bottomrule
\end{tabular}
\end{table}

Top-5 sensors hit by each removal: SAE prototype $\{0, 29, 28, 1, 30\}$;
ICA best $\{0, 28, 29, 1, 30\}$ (5/5 overlap); PCA best dispersed
(intervention is a near no-op).

So on PD31, FastICA finds a direction with the same five-sensor footprint as
the SAE prototype and damages reconstruction \emph{more} than the SAE
prototype itself ($+0.0172$ vs $+0.0117$). PCA finds a direction with the same
magnitude of geometric overlap ($|\mathrm{cos}| = 0.716$) but it is not on a
reconstruction-relevant axis.

\subsection{MEG-MMI source 2 (20 segments, original 2-subject SAE)}

\begin{table}[h]
\centering
\small
\begin{tabular}{lcc}
\toprule
\textbf{method} & \textbf{cos.\ to SAE prototype} & \textbf{mean \dmse{}} \\
\midrule
SAE prototype                  & $1.000$  & $+0.0120$ \\
best PCA component (PC25)      & $+0.721$ & $\mathbf{+0.0214}$ \\
best FastICA component (IC13)  & $-0.189$ & $+0.0056$ \\
\bottomrule
\end{tabular}
\end{table}

Top-5 sensors: SAE prototype $\{26, 27, 31, 25, 122\}$; ICA best
$\{26, 27, 25, 122, 30\}$ (4/5 overlap); PCA best
$\{27, 122, 115, 26, 31\}$ (4/5 overlap; \emph{includes sensor 115}, the
original raw source-2 ablation footprint sensor).

The MEG result inverts the PD31 pattern: PCA's most-aligned component damages
reconstruction $2\times$ the SAE prototype (and includes the external-damage
sensor 115 in its footprint), while FastICA cannot find a high-cosine direction.

\subsection{The MEG inversion is a 2-subject artifact (expanded6 update)}

Re-running the same baseline on the 6-subject expanded MEG set:

\begin{table}[h]
\centering
\small
\begin{tabular}{lcc}
\toprule
\textbf{method} & \textbf{2-subj cos.} & \textbf{expanded6 cos.} \\
\midrule
best PCA component       & $+0.721$ & $+0.454$ \\
best FastICA component   & $-0.189$ & $\mathbf{-0.748}$ \\
\bottomrule
\end{tabular}
\end{table}

With more subjects, \textbf{FastICA recovers the SAE prototype direction on
MEG} ($|\mathrm{cos}|$ jumps from $0.189$ to $0.748$) while PCA's best
alignment drops. The 2-subject ``PCA wins, ICA fails'' inversion was a
small-sample artifact, not a real EEG-vs-MEG distinction. ($\Delta$MSE
intervention values for expanded6 are unavailable due to a CPU-time failure
during the 60-min intervention loop.)

The clean cross-modality story therefore is: SAE prototype directions are
\emph{recoverable} by FastICA on standardized source-conditioned activations,
given enough cross-subject diversity; PCA is \emph{not} a uniformly reliable
substitute. The SAE pipeline is the only single method that recovers a
causally important source-prototype direction on both modalities with
comparable damage profile, both at small and at large sample size.

\begin{figure}[t]
\centering
\includegraphics[width=0.85\linewidth]{fig_pca_ica_baseline.pdf}
\caption{\textbf{SAE prototype vs PCA/ICA baselines, best-matched component.}
$\Delta$ input MSE under \emph{remove\_projection} along the SAE prototype
direction (left bars) vs the best-aligned PCA component (middle bars) vs the
best-aligned FastICA component (right bars). PD31 numbers are on the
4-subject 76-segment seed-0 SAE; MEG numbers are on the 2-subject 20-segment
seed-0 SAE. The PD31 PCA bar is slightly negative (reconstruction
\emph{improves} after PCA-component removal --- the direction is not on a
reconstruction-relevant axis).}
\label{fig:pca-ica-baseline}
\end{figure}

% =========================================================================
\section{Cross-modality contrast: spatial--temporal residual split}\label{sec:spat-temp}
% =========================================================================

BrainOmni's encoder splits its residual stream
\cite{elhage2021framework} halfway along the channel
dimension (\texttt{split\_dim} argument in the upstream code) --- the lower
half is conventionally read as ``temporal'' and the upper half as ``spatial''.
To probe whether the encoded information factorizes that way, we patch one
half of the residual stream from a clean run into a perturbed run, per
transformer block, under two perturbations: time reversal and spatial channel
permutation.

PD31 EEG, all 11 patchable blocks, 8~segments balanced across 4 subjects:

\begin{table}[h]
\centering
\small
\begin{tabular}{llcc}
\toprule
\textbf{perturbation} & \textbf{patch half} & \textbf{tiny} & \textbf{base} \\
\midrule
time\_reverse      & both     & $0.278$ & $0.241$ \\
time\_reverse      & spatial  & $0.159$ & $0.124$ \\
time\_reverse      & temporal & $0.166$ & $0.179$ \\
spatial\_mismatch  & both     & $0.269$ & $0.217$ \\
spatial\_mismatch  & spatial  & $0.146$ & $0.111$ \\
spatial\_mismatch  & temporal & $0.140$ & $0.143$ \\
\bottomrule
\end{tabular}
\caption{Mean similarity-gain by patch-half on PD31 EEG, both BrainOmni
variants.\label{tab:st-pd31}}
\end{table}

MEG-MMI, uncropped (full 30~s) 6-subject, 4 of 11 blocks:

\begin{table}[h]
\centering
\small
\begin{tabular}{llc}
\toprule
\textbf{perturbation} & \textbf{patch half} & \textbf{mean similarity gain} \\
\midrule
spatial\_mismatch  & both     & $0.209$ \\
spatial\_mismatch  & spatial  & $0.130$ \\
spatial\_mismatch  & temporal & $0.084$ \\
time\_reverse      & both     & $0.183$ \\
time\_reverse      & spatial  & $0.133$ \\
time\_reverse      & temporal & $0.082$ \\
\bottomrule
\end{tabular}
\caption{Mean similarity-gain by patch-half on uncropped 6-subject MEG-MMI.
\label{tab:st-meg}}
\end{table}

Three observations:
\begin{enumerate}[leftmargin=1.4em,itemsep=0.2em,topsep=0.3em]
\item \textbf{The ``both halves'' patch is consistently the strongest} in both
modalities and both BrainOmni variants --- encoded information is split
across the two halves rather than concentrated in one.
\item \textbf{PD31 has near-symmetric halves at \texttt{tiny} and asymmetric
halves at \texttt{base}.} At \texttt{tiny}, spatial and temporal halves are
within $0.01$ of each other for time-reverse recovery
($0.159$ vs $0.166$). At \texttt{base}, the temporal half carries more
time-reversal information ($0.179$ vs $0.124$). The bigger encoder makes the
temporal/spatial split \emph{less}, not more, symmetric.
\item \textbf{On MEG the spatial half dominates by $1.6\times$} for both
perturbations ($0.130$/$0.133$ vs $0.084$/$0.082$). Spatial mismatch is also
the stronger perturbation on MEG (both-halves gain $0.209$ vs $0.183$),
opposite of PD31. This is a genuine MEG-vs-EEG distinction visible only at
full segment duration on the 6-subject set.
\end{enumerate}

The headline mechanistic claim from the spatial--temporal split is:
\textbf{the encoder splits its residual stream into two halves that carry
overlapping but not symmetric information}, with the asymmetry sharpening at
the bigger encoder and growing larger on the higher-channel-count modality
(MEG).

% =========================================================================
\section{Decoder routing and source--sensor mismatch}\label{sec:routing}
% =========================================================================

The source-prototype results so far locate structure in the
\emph{tokenizer-input} side of the model. To check whether the same picture
holds on the decoder side, we read the \texttt{ForwardSolution} attention path
in the decoder and summarize, for each target source: which sensors place
the largest attention weight on that source, which sensors receive the largest
pre-projection contribution norm from that source, and which decoder heads
carry those routes most strongly. The framing is the standard
mechanistic-interpretability picture of features routed through circuits
\cite{olah2020zoomin,elhage2021framework}, with feed-forward layers acting
as key-value memories that rebroadcast features to specific output
neighbourhoods \cite{geva2021ffn}.

\paragraph{EEG.} On PD31, the decoder routing footprint of source~3 is
$\{0, 1, 2, 3, 29\}$ --- $3/5$ overlap with its ablation-damage footprint and
$4/5$ overlap with its source-only-reconstruction footprint. Source~2 routes
through $\{1, 3, 5, 6, 7\}$ ($1/5$ damage overlap, $3/5$ source-only overlap).
Head resolution does not overturn the picture; both source-3 heads
($H_0$, $H_2$) stay inside the same local sensor neighborhood as the aggregate
footprint.

\paragraph{MEG.} The MEG result is the cleanest mechanistic mismatch in this
work. Source~2 damages sensors $\{115, 108, 109, 244, 27\}$ when removed and
emits isolated source-only energy around $\{25\text{--}31, 122\}$ --- but the
decoder internally routes source~2 through a stable
$\{110, 111, 117, 123, 124\}$ neighborhood with $0/5$ overlap with either
external footprint. At head resolution the source splits into four stable
head-specific sensor islands ($H_0$: $\{24, 25, 26, 29, 30\}$;
$H_1$: $\{197, 199, 200, 201, 257\}$;
$H_2$: $\{98, 106, 107, 113, 114\}$, largest contribution norm;
$H_3$: $\{110, 115, 116, 122, 123\}$). Only $H_3$ touches the external
$\{115, 122\}$ neighborhood at all, and only on $1/5$ sensors. The aggregate
$\{110\text{--}124\}$ routing footprint hides a four-head decomposition.

The mechanistic implication is that ``one source = one sensor component'' is
\emph{wrong} for the BrainTokenizer's MEG path. Source~2 is transformed
through the decoder's sensor-attention layer in a way that separates internal
routing, isolated output energy, and ablation damage into at least four
head-level routes before later decoder mixing. The aggregate input/output
sensor footprints look mismatched because they are reconstructions of
different stages of the routing pipeline.

% =========================================================================
\section{Source composition: pairs and small bases}\label{sec:composition}
% =========================================================================

Sources do not act in isolation: when only the dominant pair or triad is kept,
the tokenizer reconstructs a fraction of the input proportional to that
group's energy. Across all 76 PD31 EEG segments:
\begin{itemize}
\item Pair $\{3,2\}$: pair energy fraction $40.5\,\%$, cosine to full-tokenizer
reconstruction $0.479$, cosine to the sum of single-source reconstructions
$0.945$, mean superadditive energy $-16.8\,\%$.
\item Triad $\{3, 2, 8\}$: energy fraction $54.5\,\%$, cosine to full $0.572$,
cosine to summed singles $0.892$, gain over best single $0.142$.
\item Five-source $\{3, 2, 8, 4, 13\}$: energy fraction $60.1\,\%$, cosine to
full $0.642$, cosine to summed singles $0.810$, gain over $\{3,2,8\}$ triad
$0.070$.
\end{itemize}
Across all 20 MEG-MMI segments, the pair $\{2, 11\}$ has energy fraction
$22.3\,\%$, cosine to full $0.362$, cosine to summed singles $0.940$, mean
superadditive energy $-7.5\,\%$.

Two qualitative results. \emph{(i)}~Sources combine close to additively
inside the decoder (group cosine to summed singles $\sim 0.89$--$0.94$), but
re-quantization compresses amplitude (negative superadditive energy on every
group). \emph{(ii)}~Adding sources past the dominant pair keeps closing the
reconstruction gap on EEG: $\{3,2,8\}$ beats $\{3,2\}$, and the four-source
addition that helps most is source~$13$, not source~$4$. The basis has not
plateaued at three sources.

% =========================================================================
\section{Tokenizer codebook utilization}\label{sec:codebook}
% =========================================================================

Earlier 2-subject MEG runs flagged $\sim 80\,\%$ codebook utilization
($397/430/434/442$ of $512$ codes per layer). With the 6-subject uncropped
set this rises to $489/512$, $512/512$, $509/512$, $512/512$
($95.5\,\%/100\,\%/99.4\,\%/100\,\%$). The BrainTokenizer's codebooks are
essentially fully utilized once the input distribution is broad enough; the
earlier under-utilization was a 2-subject sample artifact.

% =========================================================================
\section{Sources at the dataset level}\label{sec:dataset}
% =========================================================================

The per-source row counts in the SAE export bundles (proxy for tokenizer
``use'' of each source) are uniform by design --- the BrainTokenizer's source
allocation is balanced across rows. Figure~\ref{fig:source-dominance} shows
three complementary pictures: (a) the per-segment top-1 source counts on the
4-subject PD31 sweep, where source~3 dominates with $40/76$ wins;
(b) the per-segment top-1 source counts on the 2-subject MEG-MMI sweep, where
source~2 dominates $20/20$ (with only the top-4 sources tested); and
(c) source-ablation $\Delta$MSE on the uncropped 6-subject MEG-MMI set, where
source~9 takes over as the most-causal latent ($\Delta$MSE $0.0503$ vs
source~2's $0.0396$), with sources $\{14, 8, 15, 7, 0\}$ ranking next.
The shift from (b) to (c) is precisely the small-sample-to-expanded transition
that motivates the source-2-vs-source-9 dissociation in
Section~\ref{sec:dissociation}.

\begin{figure}[t]
\centering
\includegraphics[width=\linewidth]{fig_source_dominance.pdf}
\caption{\textbf{Tokenizer source dominance across datasets and sample sizes.}
(a)~PD31~EEG: per-segment top-1 source count, 76~segments / 4~subjects ---
source~3 dominates. (b)~MEG-MMI 2-subject 20-segment sweep: source~2 dominates
$20/20$ (only the top-4 sources by feature energy were tested per segment).
(c)~MEG-MMI uncropped 6-subject source-ablation $\Delta$MSE: source~9 emerges
as the most-causal latent at the larger sample size, with source~2 ranked
third. Source~1's negative bar reflects a small reconstruction
\emph{improvement} when source~1 is ablated (uncommon but real on this
sample).}
\label{fig:source-dominance}
\end{figure}

% =========================================================================
\section{Limitations}\label{sec:limitations}
% =========================================================================

This paper's claims rest on a small number of subjects per modality
(10 PD31, 6 MEG-MMI). Several caveats:
\begin{itemize}
\item \textbf{LOSO is brutal at $N=10$.} The PD31 leave-one-subject-out probe
collapses below chance because per-subject activation bias dominates with
only $9\,\%$ of segments held out. We rely on stratified-5-fold
(within-subject leakage) for the 1-d sufficiency claim, which makes the
\emph{bound} result (``the prototype direction carries all the diagnostic
structure the linear classifier finds'') meaningful but not a cross-subject
generalization claim.
\item \textbf{MEG LOSO at $N=6$ is meaningful.} With 6 subjects and $\sim 26$
segments each, the LOSO baseline reaches AUC~$=0.715$ (well above chance),
and the prototype-only result captures $87.7\,\%$ of that.
\item \textbf{Only \texttt{tiny} BrainOmni is fully causally analyzed.} The
\texttt{base} variant was used only for the spatial--temporal split
replication, not for prototype/patching/probe analyses. A scaling check on
\texttt{base} for PD31 plus MEG would substantially strengthen all four
mechanistic claims.
\item \textbf{The MEG 11-block spatial--temporal sweep, expanded6 PCA/ICA
$\Delta$MSE intervention, and same/diff-label MEG patching are incomplete.}
They were started during this analysis pass but killed mid-run because of
CPU contention with higher-priority items, or aborted because of script-level
path-resolution issues. Their absence is documented in the execution log;
resuming them requires a few hours of compute on a faster machine.
\item \textbf{Source-9 seed 3 is a flagged outlier.} The seed-3 expanded6
source-9 prototype has the same magnitude as the other seeds but is oriented
away from the reconstruction error axis ($\Delta$MSE $0.00078$ vs the
others' $\sim 0.025$). This is the only seed-instability we observed; the
right interpretation is probably that source~9's prototype direction is
sensitive to SAE training noise in a way source~2's is not, which itself is
informative.
\item \textbf{No comparison to subject-stratified neuroscience baselines}
(band-power probes, event-related-potential templates). The PCA/ICA
baselines are statistical decompositions on the SAE's input space, not
neuroscience-grounded controls. A clinical-linguistics audience would expect
more.
\item \textbf{Both diagnostic-label datasets are small and not balanced for
age, sex, or comorbidity.} The MEG-MMI set in particular has young
subjects (14--18 years old) which is unusual for an MDD cohort. The
diagnostic-axis interpretations should be read as ``the linear classifier
finds \emph{some} structure correlated with the binary label'' rather than
as clinical claims.
\end{itemize}

% =========================================================================
\section{Discussion and conclusion}\label{sec:discussion}
% =========================================================================

We characterized the BrainOmni tokenizer's 16-source basis through the lens of
a single \textbf{source-prototype direction} per source --- the mean of decoded
SAE features \cite{bricken2023monosemanticity,cunningham2024sae,gao2024scaling}
in the standardized activation space. Across two datasets and four SAE seeds,
this direction is seed-stable, causally necessary for input reconstruction,
diagnostic-aligned (sufficient on PD31, partially so on MEG), and
cross-subject portable through the prototype-orthogonal residual. The
direction is recoverable by FastICA on standardized activations given
sufficient cross-subject diversity, but not by PCA. Within MEG, two source
latents (sources~2 and~9) factorize qualitatively different regularities ---
subject-specific deviation vs shared population baseline --- which is the
strongest single argument that the BrainTokenizer's source basis is a real
mechanistic decomposition rather than a hash.

Two methodological takeaways for mechanistic interpretability on brain
foundation models. \emph{First}, source-level prototype directions are far
more seed-stable than individual SAE features (cosine $\geq 0.998$ vs
$\sim 0.5$ mean active mutual cosine), making them the right unit for
cross-seed mechanistic claims on these models --- consistent with the
broader observation that interpretable features are easier to reason
about as low-dimensional directions than as individual neurons
\cite{olah2020zoomin,park2024linear,marks2024sfc}. \emph{Second},
\textbf{prototype direction necessity} ($\Delta$ input MSE) and
\textbf{prototype direction sufficiency} (linear classifier 1-d AUC) are
\emph{separate} properties of the same direction; on MEG they dissociate
cleanly across two sources. Future mechanistic-interpretability work on
tokenizer-based models should report both.

\paragraph{Relation to prior brain-signal modelling.} Compared to clinical
{EEG} pipelines based on traditional decompositions
\cite{eegnet2018,schirrmeister2017} and brain-age benchmarks built on
preprocessed M/EEG \cite{engemann2022}, our analysis is the first, to our
knowledge, to ask the source-by-source question on a foundation-model
tokenizer that fuses {EEG} and {MEG}. Compared to {LaBraM}
\cite{labram2024}, {BIOT} \cite{biot2023}, {BENDR} \cite{bendr2021}, and
{Brant}/{Brant-2} \cite{brant2023,brant2_2024}, BrainOmni's RVQ
front-end is the only one that exposes a discrete latent decomposition
amenable to the prototype intervention we study here.

The headline single-axis interpretability result is the PD31 1-d projection
AUC of $0.760$ (matching the full 256-d activation classifier's $0.748$),
seed-stable across 4 SAE fits to $0.001$ AUC. We are aware of no analogous
single-axis result in the published mechanistic-interpretability-on-brain-
foundation-models literature.

We close with a refined characterization of the residual: the
prototype-orthogonal residual carries cross-subject transferable structure on
PD31 and MEG-source-2 but not on MEG-source-9; on PD31, $\sim 44\,\%$ of the
residual donor delta lives in a 1-d label-aligned axis but $\sim 87\,\%$ lives
in the orthogonal complement (subject-specific noise). The ``label-structured
residual'' claim from earlier exploratory work holds, but the bulk of the
residual donor transport is label-orthogonal.

% =========================================================================
\section*{Reproducibility}\label{sec:repro}
% =========================================================================
\addcontentsline{toc}{section}{Reproducibility}

All code is in the repository under \texttt{scripts/}. The full analysis is
reproducible end-to-end with the commands listed in the project README. The
seven figures of this paper are produced by a single script:
\begin{verbatim}
conda run -p $ENV_PATH python scripts/make_paper_figures.py --figures all
\end{verbatim}
Outputs land under \texttt{results/figures/} (PNG and PDF, 300~dpi). The four
internal reports under \texttt{reports/} document every analytic step
end-to-end:
\begin{itemize}
\item \texttt{initial\_findings.md} -- demo-sample first pass.
\item \texttt{real\_dataset\_findings.md} -- PD31 / MEG-MMI source profiles,
  pair / triad interactions, decoder routing.
\item \texttt{multi\_seed\_baseline\_replication.md} -- 4-seed SAE
  replication, PCA/ICA baselines, base-variant spatial--temporal split,
  expanded-subset analyses.
\item \texttt{research\_plan.md} -- the operational plan that produced these
  results.
\end{itemize}

% =========================================================================
\begin{thebibliography}{99}\label{sec:refs}
% =========================================================================

% ---------- Brain foundation models ---------------------------------------

\bibitem{brainomni2025}
Zhang, Z., Wang, K., et~al.
\newblock {BrainOmni}: A brain foundation model for unified {EEG} and
{MEG} signals.
\newblock In {\em Advances in Neural Information Processing Systems
(NeurIPS)}, 2025. arXiv:2505.18185.

\bibitem{labram2024}
Jiang, W.-B., Zhao, L.-M., and Lu, B.-L.
\newblock Large brain model for learning generic representations with
tremendous {EEG} data in {BCI}.
\newblock In {\em International Conference on Learning Representations
(ICLR)}, 2024 (spotlight). arXiv:2405.18765.

\bibitem{biot2023}
Yang, C., Westover, M.~B., and Sun, J.
\newblock {BIOT}: Biosignal transformer for cross-data learning in the wild.
\newblock In {\em Advances in Neural Information Processing Systems
(NeurIPS)}, 2023. arXiv:2305.10351.

\bibitem{brant2023}
Zhang, D., Yuan, Z., Yang, Y., Chen, J., Wang, J., and Li, Y.
\newblock {Brant}: Foundation model for intracranial neural signal.
\newblock In {\em Advances in Neural Information Processing Systems
(NeurIPS)}, 2023.

\bibitem{brant2_2024}
Yuan, Z., Zhang, D., Chen, J., Gu, G., and Yang, Y.
\newblock {Brant-2}: Foundation model for brain signals.
\newblock {\em arXiv:2402.10251}, 2024.

\bibitem{bendr2021}
Kostas, D., Aroca-Ouellette, S., and Rudzicz, F.
\newblock {BENDR}: Using transformers and a contrastive self-supervised
learning task to learn from massive amounts of {EEG} data.
\newblock {\em Frontiers in Human Neuroscience}, 15:653659, 2021.

\bibitem{brainlm2024}
Caro, J.~O., Ortega Caro, J., Fonseca, A.~H.~O., et~al.
\newblock {BrainLM}: A foundation model for brain activity recordings.
\newblock In {\em International Conference on Learning Representations
(ICLR)}, 2024.

\bibitem{banville2021}
Banville, H., Chehab, O., Hyv\"arinen, A., Engemann, D.-A., and Gramfort, A.
\newblock Uncovering the structure of clinical {EEG} signals with
self-supervised learning.
\newblock {\em Journal of Neural Engineering}, 18(4):046020, 2021.

\bibitem{defossez2023brain}
D\'efossez, A., Caucheteux, C., Rapin, J., Kabeli, O., and King, J.-R.
\newblock Decoding speech perception from non-invasive brain recordings.
\newblock {\em Nature Machine Intelligence}, 5:1097--1107, 2023.

\bibitem{engemann2022}
Engemann, D.~A., Mellot, A., H\"ochenberger, R., Banville, H., Sabbagh, D.,
Gemein, L., Ball, T., and Gramfort, A.
\newblock A reusable benchmark of brain-age prediction from {M/EEG}
resting-state signals.
\newblock {\em NeuroImage}, 262:119521, 2022.

\bibitem{eegnet2018}
Lawhern, V.~J., Solon, A.~J., Waytowich, N.~R., Gordon, S.~M., Hung, C.~P.,
and Lance, B.~J.
\newblock {EEGNet}: A compact convolutional neural network for {EEG}-based
brain-computer interfaces.
\newblock {\em Journal of Neural Engineering}, 15(5):056013, 2018.

\bibitem{schirrmeister2017}
Schirrmeister, R.~T., Springenberg, J.~T., Fiederer, L. D.~J., et~al.
\newblock Deep learning with convolutional neural networks for {EEG}
decoding and visualization.
\newblock {\em Human Brain Mapping}, 38(11):5391--5420, 2017.

% ---------- Datasets and clinical ----------------------------------------

\bibitem{ds002778}
Rockhill, A.~P., Jackson, N., George, J.~S., Aron, A.~R., and Swann, N.~C.
\newblock {UC} San Diego resting-state {EEG} data from patients with
{Parkinson}'s disease.
\newblock {\em OpenNeuro}, dataset \texttt{ds002778}, 2018.

\bibitem{ds003568}
Liuzzi, L., Pine, D.~S., Fox, N.~A., and Furman, D.~J.
\newblock Mood induction in {MDD} and healthy adolescents
(\texttt{task-mmi3} {MEG} dataset).
\newblock {\em OpenNeuro}, dataset \texttt{ds003568}, 2021.

\bibitem{knutson2001mid}
Knutson, B., Adams, C.~M., Fong, G.~W., and Hommer, D.
\newblock Anticipation of increasing monetary reward selectively recruits
nucleus accumbens.
\newblock {\em Journal of Neuroscience}, 21(16):RC159, 2001.

% ---------- Sparse coding, ICA, classical decomposition -------------------

\bibitem{olshausen1996}
Olshausen, B.~A. and Field, D.~J.
\newblock Emergence of simple-cell receptive field properties by learning a
sparse code for natural images.
\newblock {\em Nature}, 381(6583):607--609, 1996.

\bibitem{bell1995infomax}
Bell, A.~J. and Sejnowski, T.~J.
\newblock An information-maximization approach to blind separation and blind
deconvolution.
\newblock {\em Neural Computation}, 7(6):1129--1159, 1995.

\bibitem{hyvarinen2000}
Hyv\"arinen, A. and Oja, E.
\newblock Independent component analysis: Algorithms and applications.
\newblock {\em Neural Networks}, 13(4--5):411--430, 2000.

\bibitem{makeig1996}
Makeig, S., Bell, A.~J., Jung, T.-P., and Sejnowski, T.~J.
\newblock Independent component analysis of electroencephalographic data.
\newblock In {\em Advances in Neural Information Processing Systems
(NIPS)}, vol.~8, pages 145--151, 1996.

% ---------- Vector quantization & residual VQ -----------------------------

\bibitem{vqvae2017}
van~den Oord, A., Vinyals, O., and Kavukcuoglu, K.
\newblock Neural discrete representation learning.
\newblock In {\em Advances in Neural Information Processing Systems
(NIPS)}, 2017. arXiv:1711.00937.

\bibitem{vqvae2_2019}
Razavi, A., van~den Oord, A., and Vinyals, O.
\newblock Generating diverse high-fidelity images with {VQ-VAE-2}.
\newblock In {\em Advances in Neural Information Processing Systems
(NeurIPS)}, 2019. arXiv:1906.00446.

\bibitem{soundstream2021}
Zeghidour, N., Luebs, A., Omran, A., Skoglund, J., and Tagliasacchi, M.
\newblock {SoundStream}: An end-to-end neural audio codec.
\newblock {\em IEEE/ACM Transactions on Audio, Speech, and Language
Processing}, 30:495--507, 2021. arXiv:2107.03312.

\bibitem{encodec2022}
D\'efossez, A., Copet, J., Synnaeve, G., and Adi, Y.
\newblock High fidelity neural audio compression.
\newblock {\em arXiv:2210.13438}, 2022.

% ---------- Mechanistic interpretability foundations ----------------------

\bibitem{olah2020zoomin}
Olah, C., Cammarata, N., Schubert, L., Goh, G., Petrov, M., and Carter, S.
\newblock Zoom in: An introduction to circuits.
\newblock {\em Distill}, 2020.

\bibitem{elhage2021framework}
Elhage, N., Nanda, N., Olsson, C., Henighan, T., Joseph, N., et~al.
\newblock A mathematical framework for transformer circuits.
\newblock {\em Transformer Circuits Thread, Anthropic}, 2021.

% ---------- Sparse autoencoders for interpretability ----------------------

\bibitem{bricken2023monosemanticity}
Bricken, T., Templeton, A., Batson, J., et~al.
\newblock Towards monosemanticity: Decomposing language models with
dictionary learning.
\newblock {\em Transformer Circuits Thread, Anthropic}, 2023.

\bibitem{cunningham2024sae}
Cunningham, H., Ewart, A., Riggs Smith, L., Huben, R., and Sharkey, L.
\newblock Sparse autoencoders find highly interpretable features in language
models.
\newblock In {\em International Conference on Learning Representations
(ICLR)}, 2024. arXiv:2309.08600.

\bibitem{templeton2024scaling}
Templeton, A., Conerly, T., Marcus, J., Lindsey, J., et~al.
\newblock Scaling monosemanticity: Extracting interpretable features from
{Claude 3 Sonnet}.
\newblock {\em Transformer Circuits Thread, Anthropic}, 2024.

\bibitem{gao2024scaling}
Gao, L., Dupr\'e la Tour, T., Tillman, H., Goh, G., Troll, R., Radford, A.,
Sutskever, I., Leike, J., and Wu, J.
\newblock Scaling and evaluating sparse autoencoders.
\newblock {\em arXiv:2406.04093}, 2024.

\bibitem{rajamanoharan2024gated}
Rajamanoharan, S., Conmy, A., Smith, L., Lieberum, T., Varma, V., Kram\'ar, J.,
Shah, R., and Nanda, N.
\newblock Improving dictionary learning with gated sparse autoencoders.
\newblock {\em arXiv:2404.16014}, 2024.

\bibitem{lieberum2024gemma}
Lieberum, T., Rajamanoharan, S., Conmy, A., Smith, L., Sonnerat, N., Varma, V.,
Kram\'ar, J., Dragan, A., Shah, R., and Nanda, N.
\newblock {Gemma Scope}: Open sparse autoencoders everywhere all at once
on {Gemma 2}.
\newblock {\em arXiv:2408.05147}, 2024.

\bibitem{marks2024sfc}
Marks, S., Rager, C., Michaud, E.~J., Belinkov, Y., Bau, D., and Mueller, A.
\newblock Sparse feature circuits: Discovering and editing interpretable
causal graphs in language models.
\newblock In {\em International Conference on Learning Representations
(ICLR)}, 2025. arXiv:2403.19647.

\bibitem{park2024linear}
Park, K., Choe, Y.~J., and Veitch, V.
\newblock The linear representation hypothesis and the geometry of large
language models.
\newblock In {\em International Conference on Machine Learning (ICML)},
2024. arXiv:2311.03658.

\bibitem{geva2021ffn}
Geva, M., Schuster, R., Berant, J., and Levy, O.
\newblock Transformer feed-forward layers are key-value memories.
\newblock In {\em Conference on Empirical Methods in Natural Language
Processing (EMNLP)}, 2021. arXiv:2012.14913.

% ---------- Causal interventions, activation patching --------------------

\bibitem{rome2022}
Meng, K., Bau, D., Andonian, A., and Belinkov, Y.
\newblock Locating and editing factual associations in {GPT}.
\newblock In {\em Advances in Neural Information Processing Systems
(NeurIPS)}, 2022. arXiv:2202.05262.

\bibitem{conmy2023acdc}
Conmy, A., Mavor-Parker, A.~N., Lynch, A., Heimersheim, S., and
Garriga-Alonso, A.
\newblock Towards automated circuit discovery for mechanistic
interpretability.
\newblock In {\em Advances in Neural Information Processing Systems
(NeurIPS)}, 2023 (spotlight). arXiv:2304.14997.

\bibitem{wang2023ioi}
Wang, K.~R., Variengien, A., Conmy, A., Shlegeris, B., and Steinhardt, J.
\newblock Interpretability in the wild: A circuit for indirect object
identification in {GPT-2 small}.
\newblock In {\em International Conference on Learning Representations
(ICLR)}, 2023. arXiv:2211.00593.

\bibitem{vig2020causal}
Vig, J., Gehrmann, S., Belinkov, Y., Qian, S., Nevo, D., Singer, Y., and
Shieber, S.~M.
\newblock Investigating gender bias in language models using causal
mediation analysis.
\newblock In {\em Advances in Neural Information Processing Systems
(NeurIPS)}, 2020. arXiv:2004.12265.

\bibitem{geiger2024abstraction}
Geiger, A., Ibeling, D., Zur, A., Chaudhary, M., Chauhan, S., Huang, J.,
Arora, A., Wu, Z., Goodman, N., Potts, C., and Icard, T.
\newblock Causal abstraction: A theoretical foundation for mechanistic
interpretability.
\newblock {\em Journal of Machine Learning Research}, 26:1--65, 2025.
arXiv:2301.04709.

% ---------- Probing classifiers ------------------------------------------

\bibitem{alain2017probes}
Alain, G. and Bengio, Y.
\newblock Understanding intermediate layers using linear classifier probes.
\newblock In {\em ICLR Workshop}, 2017. arXiv:1610.01644.

\bibitem{belinkov2022probing}
Belinkov, Y.
\newblock Probing classifiers: Promises, shortcomings, and advances.
\newblock {\em Computational Linguistics}, 48(1):207--219, 2022.

\bibitem{hewitt2019probes}
Hewitt, J. and Liang, P.
\newblock Designing and interpreting probes with control tasks.
\newblock In {\em Conference on Empirical Methods in Natural Language
Processing (EMNLP-IJCNLP)}, pages 2733--2743, 2019.
arXiv:1909.03368.

\end{thebibliography}

\end{document}
